\chapter{Radon representation from positive functionals}\label{sec:radon-representation}
The examples of measures become useful in compactness arguments because one can construct a measure from limits of its integrals. We prove the positive Riesz representation theorem on an open set $U\subset\mathbb R^n$, following \cite{Sim18}, Chapter 1, Section 5, pp.~30--37. The elementary Caratheodory construction of a measure on the measurable sets of an outer measure is a background input. Metric separation and all regularity and representation steps used here are proved below.

\begin{definition}[A positive functional on compactly supported tests]\label{def:positive-test-functional}
A positive linear functional is a map $L:C_c(U;\mathbb R)\to\mathbb R$ which is linear and satisfies $L(f)\geq0$ whenever $f\geq0$. Values are finite real numbers. No single bound $|L(f)|\leq C\|f\|_\infty$ for all supports is assumed. Positivity will instead give bounds on each fixed compact support. See \cite{Sim18}, Theorem 5.12 and Remark 5.13.
\end{definition}

\begin{lemma}[Compact cutoffs and a finite partition of unity]\label{lem:radon-compact-cutoffs}
If $K\subset U$ is compact and $O\subset U$ is open with $K\subset O$, there is $\eta\in C_c(O)$, $0\leq\eta\leq1$, equal to one on a neighborhood of $K$. If finitely many open sets $O_1,\ldots,O_N$ cover $K$, there are $\eta_i\in C_c(O_i)$, $0\leq\eta_i\leq1$, whose sum is one on a neighborhood of $K$ and is at most one everywhere.
\end{lemma}
\begin{proof}
The empty-set cases are immediate. For the first assertion choose $\delta>0$ so that the closed $\delta$-neighborhood of $K$ is compact and lies in $O$. A continuous piecewise linear function of $\operatorname{dist}(x,K)$, equal to one for distance at most $\delta/2$ and zero for distance at least $\delta$, is the required cutoff. For the second assertion, cover $K$ by finitely many small balls whose closures each lie in one of the $O_i$. Grouping their cutoffs gives nonnegative $\psi_i\in C_c(O_i)$ with $S=\sum_i\psi_i>0$ on a neighborhood of $K$. Apply the first assertion to obtain $\chi\in C_c(\{S>0\})$, equal to one near $K$ and between zero and one. Set $\eta_i=\chi\psi_i/S$ where $S>0$, and zero elsewhere. Its extension is continuous because $\operatorname{spt}\chi$ stays in $\{S>0\}$. The sum equals $\chi$. This is the Euclidean version of \cite{Sim18}, Lemma 5.4 and Corollary 5.5.
\end{proof}

\begin{lemma}[Positivity gives continuity on a fixed support]\label{lem:positive-functional-local-bound}
If $L$ is positive, $K\subset U$ is compact, and $\eta\in C_c(U)$ is nonnegative and at least one on $K$, then
\[
 |L(f)|\leq L(\eta)\|f\|_\infty\quad\text{for every }f\in C_c(U)\text{ supported in }K.
\]
Thus uniform convergence with a common compact support implies convergence of the functional values.
\end{lemma}
\begin{proof}
The pointwise inequalities $-\|f\|_\infty\eta\leq f\leq\|f\|_\infty\eta$ and positivity give the bound. Apply it to the difference of two functions for the continuity assertion. This is \cite{Sim18}, Remark 5.13.
\end{proof}

\begin{proposition}[The open-set mass associated with a positive functional]\label{prop:functional-open-mass}
For open $O\subset U$, put
\[
 m(O)=\sup\{L(f):f\in C_c(O),\ 0\leq f\leq1\}.
\]
Then $m$ is monotone, $m(\varnothing)=0$, is countably subadditive on open sets, and is additive on disjoint open sets, allowing infinite values. The function
\[
 \mu^*(A)=\inf\{m(O):A\subset O\subset U,\ O\text{ open}\}
\]
for arbitrary $A\subset U$ is an outer measure, equals $m$ on open sets, is finite on compact sets, and satisfies
\[
 \mu^*(A\cup B)=\mu^*(A)+\mu^*(B)
 \quad\text{if }\operatorname{dist}(A,B)>0.
\]
\end{proposition}
\begin{proof}
Monotonicity and the empty-set assertion are immediate. If $O\subset\bigcup_jO_j$ and $f$ is admissible for $m(O)$, compactness of its support and Lemma \ref{lem:radon-compact-cutoffs} give a finite subordinate partition, with $f=\sum_j\eta_jf$. Therefore $L(f)\leq\sum_jm(O_j)$; take the supremum. For disjoint opens, finitely many admissible tests have sum at most one, so taking their separate suprema gives $m(\bigcup_jO_j)\geq\sum_{j=1}^Nm(O_j)$ for every finite $N$. This proves countable additivity on disjoint opens.

The infimum definition is consistent on opens by monotonicity. To prove countable subadditivity of $\mu^*$, the only nontrivial case has $\sum_j\mu^*(A_j)<\infty$. Choose $O_j\supset A_j$ with $m(O_j)\leq\mu^*(A_j)+\varepsilon2^{-j}$, use open-set subadditivity, and let $\varepsilon\downarrow0$. For compact $K$, choose an open neighborhood $W$ with compact closure in $U$ and a cutoff $\eta\geq1$ on $\overline W$. Every test for $m(W)$ is at most $\eta$, so $\mu^*(K)\leq m(W)\leq L(\eta)<\infty$. Finally, if $A,B$ have positive distance, choose disjoint open neighborhoods $V_A,V_B$ of them. For every open $O\supset A\cup B$,
\[
 m(O)\geq m(O\cap V_A)+m(O\cap V_B)\geq\mu^*(A)+\mu^*(B).
\]
Taking the infimum and using subadditivity proves the metric separation identity. This follows the construction in \cite{Sim18}, proof of Theorem 5.12.
\end{proof}

\begin{lemma}[Metric outer measures make Borel sets measurable]\label{lem:metric-outer-borel}
Let $\mu^*$ be an outer measure on a metric space, additive on pairs of sets at positive distance. Then every closed set, and hence every Borel set, is Caratheodory measurable.
\end{lemma}
\begin{proof}
Fix a closed set $F$ and a set $E$ with $\mu^*(E)<\infty$. The empty and whole-space cases are immediate, so assume $F$ is nonempty and proper. Write $d(x)=\operatorname{dist}(x,F)$ and
\[
 A_j=E\cap\{1/(j+1)<d\leq1/j\},\qquad j\geq1.
\]
Any finite family of these annuli with indices of the same parity consists of sets separated from one another by positive distances, because $d$ is 1-Lipschitz. Their outer measures sum to at most $\mu^*(E)$. Therefore $\sum_{j\geq1}\mu^*(A_j)<\infty$, and the tails tend to zero. The part of $E\setminus F$ with $d>1/k$ is separated from $E\cap F$; hence
\[
 \mu^*(E)\geq\mu^*(E\cap F)+\mu^*(E\cap\{d>1/k\}).
\]
The latter set together with $\bigcup_{j\geq k}A_j$ covers $E\setminus F$. Subadditivity and the vanishing tail now give $\mu^*(E)\geq\mu^*(E\cap F)+\mu^*(E\setminus F)$. The reverse inequality is subadditivity. For $\mu^*(E)=\infty$ the required greater-than-or-equal inequality is automatic. Thus $F$ is measurable. The Caratheodory measurable sets form a sigma algebra, so all Borel sets are measurable. Compare \cite{Sim18}, Chapter 1, Theorem 1.15.
\end{proof}

\begin{proposition}[The constructed measure is Radon]\label{prop:functional-measure-radon}
The restriction of $\mu^*$ from Proposition \ref{prop:functional-open-mass} to Borel sets is a Radon measure $\mu$. In particular, for every open $O\subset U$,
\[
 \mu(O)=\sup_{K\subset O,\ K\text{ compact}}\mu(K)=m(O).
\]
\end{proposition}
\begin{proof}
Lemma \ref{lem:metric-outer-borel} gives a Borel measure, and Proposition \ref{prop:functional-open-mass} gives local finiteness and outer regularity. If $f\in C_c(O)$, $0\leq f\leq1$, with compact support $K$, then $L(f)\leq m(V)$ for every open $V\supset K$, so $L(f)\leq\mu(K)$. Taking the supremum in $f$ proves inner regularity on open sets; the opposite inequality is monotonicity.

For a Borel set $A$ of finite measure, choose open $O\supset A$ with $\mu(O\setminus A)<\varepsilon$, compact $K\subset O$ with $\mu(O\setminus K)<\varepsilon$, and open $V\supset O\setminus A$ with $\mu(V)<2\varepsilon$. Then $K\setminus V$ is compact and contained in $A$, and $\mu(A\setminus(K\setminus V))<3\varepsilon$. Thus $A$ is inner regular. For arbitrary Borel $A$, exhaust $U$ by compact sets $K_j$. The sets $A\cap K_j$ have finite measure and increase to $A$, so their inner regularity and continuity from below give the general assertion. Compare \cite{Sim18}, Lemmas 5.7--5.9.
\end{proof}

\begin{theorem}[Positive Riesz representation on a Euclidean open set]\label{thm:positive-riesz-representation}
For every positive linear functional $L:C_c(U;\mathbb R)\to\mathbb R$ there is a unique Radon measure $\mu$ on $U$ such that
\[
 L(f)=\int_U f\,d\mu\qquad(f\in C_c(U;\mathbb R)).
\]
Its values on open sets are given by $\mu(O)=\sup\{L(\phi):\phi\in C_c(O),\ 0\leq\phi\leq1\}$. This is \cite{Sim18}, Theorem 5.12, specialized to a Euclidean open set.
\end{theorem}
\begin{proof}
Construct the Radon measure by Proposition \ref{prop:functional-measure-radon}. First note that, for $0\leq h\leq1$ with compact support,
\[
 L(h)\leq\mu(\{h>0\}).
\]
Indeed $(h-\varepsilon)_+$ has compact support in the open set $\{h>0\}$ and is bounded by one, while the difference from $h$ tends uniformly to zero on a common compact support. Lemma \ref{lem:positive-functional-local-bound} gives the inequality upon taking the limit.

For $f\geq0$ with compact support, fix $\delta>0$ and an integer $N$ with $N\delta\geq\max f$. Set $h_i=\min(1,(f-i\delta)_+/\delta)$ for $0\leq i<N$. Pointwise $f=\delta\sum_i h_i$. The upper bound just proved and the fact that $h_i=1$ on $\{f>(i+1)\delta\}$ give
\[
 \delta\sum_{i=0}^{N-1}\mu(\{f>(i+1)\delta\})
 \leq L(f)\leq
 \delta\sum_{i=0}^{N-1}\mu(\{f>i\delta\}).
\]
For the lower bound, every admissible test supported in that open superlevel set is at most $h_i$. The same two sums bound $\int f\,d\mu$ by pointwise simple-function comparison. Their difference is at most $\delta\mu(\{f>0\})$, which tends to zero by local finiteness. Thus $L(f)=\int f\,d\mu$. Subtract positive and negative parts for signed $f$.

Finally any representing Radon measure satisfies the displayed open-set formula: the upper bound is immediate, and cutoffs equal to one on compact subsets yield the reverse bound by inner regularity. Equal open-set masses and outer regularity imply equal Borel masses. This proves uniqueness.
\end{proof}

\begin{theorem}[Continuous compactly supported tests are dense in finite-power spaces]\label{thm:radon-lp-density}
For any Radon measure $\mu$ on $U$ and $1\leq p<\infty$, $C_c(U)$ is dense in $L^p(\mu)$. The assertion concerns the completed measure as well, with functions identified almost everywhere. See \cite{Sim18}, Theorem 5.10.
\end{theorem}
\begin{proof}
Truncation of values and support, followed by approximation by simple functions, reduces the claim to $\mathbf1_A$ with $\mu(A)<\infty$; a completed-measurable set can first be replaced by a Borel set modulo a null set. Choose compact $K\subset A$ and open $O\supset A$ such that $\mu(A\setminus K)<\varepsilon$ and $\mu(O\setminus A)<\varepsilon$. Lemma \ref{lem:radon-compact-cutoffs} gives $\phi\in C_c(O)$, $0\leq\phi\leq1$, equal to one on $K$. Then
\[
 \int|\mathbf1_A-\phi|^p\,d\mu\leq\mu(O\setminus K)<2\varepsilon.
\]
Finite linear combinations approximate simple functions by the $L^p$ triangle inequality. This proves the result. The finite-power assumption is essential to this approximation argument; no density in $L^\infty$ is asserted.
\end{proof}
